\section{All orders in the directed subtree contraction}\label{me:directed}
We recall exactly what the directed certificate controls. Put
$a_0=1/(2\pi)$, $k_0=2^{-16}$, $R_0=1/8$, $b_0=2^{-11}$,
$r_0=1/16$, $t_0=2^{-14}$ and $\chi=1-\cos r_0$.
Choose a root and, at every nonroot vertex, $0<c_v<k_0$ and
$\delta_v=\min\{b_0/2,c_v\chi/[4(1+c_v)]\}$. Let $q_v$ be
one at a nonroot vertex and zero at the root. With children
$\operatorname{Ch}(v)$ define
\begin{align}
 a_v&=\frac{d_v}{2\pi},& b_v&=\frac{n_v}{2\pi},&D_v&=a_v-\ii b_v,
       \notag\\
 A_v&=\frac2{\pi^2}\left(q_vk_0+\frac12\sum_{w\in\operatorname{Ch}(v)}c_w\right),&
 B_v&=16q_v+|\operatorname{Ch}(v)|,&
 C_v&=B_v+\frac{b_v^2}{4A_v}.\label{me:directed-data}
\end{align}
All $A_v,B_v$ are positive because the tree has an edge. Put
$\delta_v^{\rm in}=\min(\{b_0\}\cup\{\delta_w:w\in\operatorname{Ch}(v)\})$,
\begin{align}
 y_v&=\min\left\{\frac{\delta_v^{\rm in}}2,
 \frac{t_0|D_v|}{4C_v},
 \frac{A_v\pi^2}{8(\pi|b_v|+t_0|D_v|)},
 \sqrt{\frac{A_v\pi^2}{8B_v}}\right\},\notag\\
 \beta_v&=\frac{210}{361}\frac{a_v}{|D_v|^2},\qquad
 \eta_v=\min\left\{\tfrac12t_0|D_v|y_v,
 \tfrac12 A_v\pi^2,A_vr_0^2,\tfrac18a_vr_0^2\right\}.
 \label{me:directed-losses}
\end{align}
The criterion is $t_0|D_v|\ge2$ and $\beta_v<\eta_v/2$ at every
vertex, and $\beta_v<(k_0-c_v)\chi$ at each nonroot vertex.
These are the same finite inequalities as in \cite{CircleTree}.

The full edge polynomial is
\begin{equation}\label{me:full-kernel}
 \mathscr K_N(t,u;s)=\sum_{k=0}^{N-1}
       e_k(-2\cos t)e_k(-2\cos u)\frac{H_N(s)}{W_k(s)}.
\end{equation}
Its trigonometric degree in either angle is at most $N-1$;
there is no high-index truncation in this definition. Let
$M_{v,N}(u;s)$ denote the exact polynomial sent by a subtree to
its parent angle $u$. It is the actual trace of the parent kernel
times the exact child messages, with the original typed scalar.
At the root the parent kernel is absent. This is an evaluation
order for a fixed contraction, not a new diagram.

The written induction in \cite[Sections 7.3--7.5]{CircleTree} proves,
on the parent rectangle $|\Re u|\le R_0$, $|\Im u|\le b_0$,
a nonzero conditional first-charge saddle. Its local real second
derivative is at most $-(19/20)a_v$ and its complex Hessian modulus
is at least $(19/20)|D_v|$. All tail, endpoint, truncation and
connector errors are below that first charge by a positive margin
at least $\eta_v-\beta_v$. Each retained charge $\ell\ge2$ has
additional loss
\begin{equation}\label{me:directed-charge-loss}
 \frac{190}{441}\frac{a_v}{|D_v|^2}(\ell^2-1).
\end{equation}
The outgoing real-envelope margin is
$(k_0-c_v)\chi-\beta_v>0$ when an outgoing edge exists. There is
no outgoing requirement at the root. These assertions are the
conditional conclusions of that proof, not consequences inferred
solely from its final scalar asymptotic. The full kernel, all
normal data and the exact child scalars are retained there.
Lemma~\ref{me:translated} verifies these same finite interfaces
on the translated parameter compacts. For one fixed finite tree
all these finitely many positive margins have positive minimum.
The infinite charge tail is one geometric series, not a minimum
of infinitely many unrelated gaps.

\begin{proposition}[Normalized induction to arbitrary order]
\label{me:directed-orders}
Every fixed directed certificate above gives all relative orders
of the actual normalized subtree messages on smaller parent
rectangles and the meridian compacts of Theorem~\ref{me:all-orders}.
Their phase is stationary elimination of \eqref{me:base-action},
with the exact scalar factors restored at the root.
\end{proposition}
\begin{proof}
Strengthen the induction to
\begin{equation}\label{me:message-orders}
 \frac{M_{v,N}(u;s)}{M_{v,N}(0;s)}
 =e^{Nh_v(u)}\frac{\alpha_v(u)}{\alpha_v(0)}
 \left(1+\sum_{j=1}^L N^{-j}a_{v,j}(u,s)
                          +O(N^{-L-1})\right).
\end{equation}
Here $h_v(0)=0$, $\alpha_v$ is nonzero, and the coefficients and
remainders are holomorphic. The exact denominator is nonzero
for all sufficiently large levels by the conditional first-charge
estimate. It is never replaced by an estimate for its modulus.

For a terminal vertex, Lemma~\ref{me:local-expansion} applied to
its full kernel and trace proves \eqref{me:message-orders}; its
complement and currents retain the original positive margins.
Suppose it holds for all children. Their finite product has an
expansion to every chosen relative order, with a nonzero normalized
leading amplitude. On the first-charge local contour use that
expansion and the canonical parent-kernel expansion. The conditional
Hessian and sine factor are nonzero. Absolute Gaussian domination
from Lemma~\ref{me:local-expansion} preserves the relative remainder,
and extra derivatives are obtained by beginning at a higher order.
This proves the first-charge expansion and its exact normalization
at $u=0$.

Outside that contour continue to use the exact child polynomials
and their established real and strip envelopes. Taylor polynomials
for the child expansion are not substituted on a tail. Nonminimal
charges are their local contributions plus additive errors already
below the first charge; a relative expansion at an exponentially
small nonminimal saddle is unnecessary. Inequality
\eqref{me:directed-charge-loss} and the original endpoint and
regeneration margins make these errors smaller than any prescribed
relative inverse power. The outgoing-envelope step again uses the
exact normalized message. It is unaffected by the additional local
Taylor coefficients. This closes the induction on the finite tree.

To identify the joint phase, restore the wire variables in each
parent kernel. Their Hessian blocks are invertible. At each angular
elimination the conditional Hessian is invertible by the displayed
bound. The block determinant identity therefore proves that the
original joint Hessian is invertible, and back-substitution recovers
its selected critical point. No equation for a fixed total action
was guessed from the local critical values. Exact factors
$M_{v,N}(0;s)$ restore every constant subtracted in the normalized
messages. At any finite order, repeated Gaussian integration gives
the same moments as the joint local integral: integrating an
invertible block Gaussian and then its Schur complement is the
usual completion of the square, and differentiation generates all
polynomial moments. The uniform absolute bounds already proved
justify that finite repeated integration. Thus both the phase and
the lower coefficients belong to the original joint source.
\end{proof}

This argument allows constants depending on the fixed tree's size
and depth, as an all-orders theorem for one fixed diagram should.
It neither strengthens the directed criterion nor applies the
summed criterion to a tree which only satisfies local inequalities.
